% Lecture 10: exceptions, error contracts, and RAII
\section{Лекция 10. Ошибки API: status, exceptions и RAII}
\label{sec:lecture-10}

До сих пор примеры в основном шли по успешному пути. Но файл может отсутствовать,
данные могут быть испорчены, а запись --- запрещена операционной системой. Нужно
решить, как сообщить об ошибке caller, провести её через несколько уровней и кто
освободит уже полученные ресурсы.

\begin{importantbox}[Причинная линия]
Успешный алгоритм $\rightarrow$ ошибка $\rightarrow$ проверки status на каждом
уровне $\rightarrow$ отделение failure path через exception $\rightarrow$ stack
unwinding $\rightarrow$ destructors локальных objects $\rightarrow$ RAII
$\rightarrow$ осознанный выбор error contract.
\end{importantbox}

Exceptions не «лучше» status codes. RAII также не требует exceptions: он
обеспечивает cleanup при normal exit, early return и unwinding. Это две разные
design decisions: как переносить failure и как владеть ресурсом.

\subsection{Ошибка как часть contract API}

Слово «ошибка» скрывает разные ситуации:
обычный допустимый результат: поиск нашёл 0 совпадений; отсутствие значения в domain semantics: лекция 11 ещё не заполнена; ожидаемое recoverable condition: неверный ввод можно повторить; programming error: caller нарушил documented precondition; exceptional failure операции: обязательный отчёт нельзя прочитать; environment failure: file missing, permission denied, malformed input.

Это не абсолютная классификация. Missing file может быть нормальным результатом
поиска optional config и failure обязательной загрузки. Поэтому способ сообщения
--- часть contract. Caller должен знать return type, возможные exceptions и
гарантии о state после неуспеха. \texttt{assert} проверяет programming invariant
в подходящих сборках; это не общий транспорт recoverable errors.

\subsection{Baseline: явный status}

В C++17 нет standard \texttt{std::expected}. Небольшому API достаточно простой
структуры, не имитирующей целую библиотеку:

\begin{cppcode}
struct Status {
    bool ok;
    std::string message;
};

Status read_config(const Path& path, std::string& value);
\end{cppcode}

Полезный result передаётся через reference, а return несёт success/failure. В
нескольких слоях появляется propagation:

\begin{cppcode}
Status transform(const Path& input, const Path& output) {
    std::string value;
    Status status = read_config(input, value);
    if (!status.ok) return status;
    // normal algorithm
    return {true, {}};
}

int command(const Path& input, const Path& output) {
    const Status status = transform(input, output);
    if (!status.ok) {
        std::cerr << status.message << '\n';
        return 2;
    }
    return 0;
}
\end{cppcode}

Плюсы: control flow виден в interface и source; частый ожидаемый исход удобно
обрабатывать рядом. Цена: checks повторяются, пропущенная проверка теряет failure,
normal value и error transport приходится совмещать. \texttt{[[nodiscard]]}
помогает заметить игнорирование результата, но не проектирует result за автора.

\subsection{Exceptions: синтаксис и semantics}

\begin{cppcode}
std::string read_required(const Path& path) {
    std::ifstream input{path};
    if (!input) {
        throw std::runtime_error{"cannot open " + path.string()};
    }
    return read_value(input);
}

try {
    build_report(path);
} catch (const std::invalid_argument& error) {
    std::cerr << "bad data: " << error.what() << '\n';
} catch (const std::exception& error) {
    std::cerr << "operation failed: " << error.what() << '\n';
}
\end{cppcode}

\texttt{throw expression} инициализирует exception object и начинает поиск
matching handler. После \texttt{throw} normal path текущего scope не продолжается.
\texttt{try} отмечает guarded statements, \texttt{catch} --- handler. Более
specific handler ставят раньше общего.

Если подходящего handler в текущей функции нет, поиск продолжается во внешнем
динамическом call chain. Промежуточный уровень обычно не ловит failure, который
не умеет исправить. Если exception достигает верхней границы без handler,
runtime вызывает механизм завершения, включающий \texttt{std::terminate}; это
точнее, чем обещание «просто crash» определённого вида.

Для standard hierarchy обычна форма
\texttt{catch (const std::exception\& error)}. Reference избегает copy и slicing,
\texttt{const} запрещает изменение, virtual \texttt{what()} даёт сообщение.
Raw integers и strings технически допустимы, но это плохой production contract.
Когда смысл совпадает, используют \texttt{std::runtime\_error},
\texttt{std::invalid\_argument}; знакомый library example
\texttt{std::out\_of\_range} бросает, например, \texttt{vector::at}.

\texttt{catch (...)} совпадает с любым exception. Он нужен редко, например для
boundary action и rethrow. Он не даёт object для \texttt{what()} и не оправдывает
молчаливое проглатывание ошибки.

\subsection{Rethrow}

\begin{cppcode}
void middle() {
    try {
        low_level();
    } catch (const std::exception& error) {
        std::cerr << "middle context: " << error.what() << '\n';
        throw;
    }
}
\end{cppcode}

Пустой \texttt{throw;} внутри handler повторно бросает текущее exception и
сохраняет его type. \texttt{throw error;} --- новый throw из expression:
object копируется по static type expression, поэтому возможно slicing. Rethrow
уместен, если слой добавил локальный контекст, но не способен восстановиться.

\subsection{Stack unwinding и lifetime}

Когда exception покидает scope, уничтожаются уже сконструированные automatic
objects, lifetime которых заканчивается из-за выхода. Их destructors вызываются
в порядке, обратном construction. Затем поиск handler продолжается наружу. Это
называется \term{stack unwinding}.

\begin{cppcode}
void low() {
    Trace low{"low"};
    throw std::runtime_error{"deep failure"};
}

void middle() {
    Trace middle{"middle"};
    low();
}
\end{cppcode}

Эксперимент печатает:
\begin{outputcode}
+ top
+ middle
+ low
- low
- middle
- top
top boundary caught: deep failure
\end{outputcode}

Это наблюдение destructors, а не требование реализовать automatic storage как
конкретные CPU stack frames. Standard задаёт lifetime и destruction; ABI и
machine representation выбирает implementation.

\subsection{RAII: ресурс принадлежит object lifetime}

\term{RAII} --- Resource Acquisition Is Initialization. Owner получает ресурс
при construction (или создаётся сразу в корректном owning state), хранит его и
освобождает в destructor. Scope задаёт lifetime. Cleanup срабатывает при обычном
выходе, early return и unwinding.

Resource --- не только heap memory. Это file handle, temporary file, mutex lock,
socket/OS handle, transaction guard или logging scope. Smart pointers --- лишь
будущий пример RAII; для понимания механизма они не нужны.

Без owner алгоритм обрастает ветками:

\begin{cppcode}
Handle h = acquire();
if (!operation_a(h)) {
    release(h);
    return failure;
}
operation_b(h);             // may throw: кто вызовет release?
release(h);
\end{cppcode}

С RAII основной алгоритм не отвечает за release:

\begin{cppcode}
ResourceOwner owner{acquire()};
if (!operation_a(owner)) return failure;
operation_b(owner);         // destructor работает и при throw
\end{cppcode}

Standard library уже использует RAII. \texttt{std::ifstream}/\texttt{ofstream}
закрывают файл, \texttt{std::vector} и \texttt{std::string} освобождают storage.
Не нужен wrapper без новой ответственности. Учебный \texttt{TemporaryReport}
добавляет policy: удалить temporary artifact, если \texttt{commit} не состоялся.

\subsection{RAII-owner и transactional thinking}

Owner из \texttt{error\_model\_lab} открывает \texttt{report.tmp} через
\texttt{std::ofstream}. Copy запрещена через \texttt{= delete}: два owners не
должны владеть одним файлом. Move semantics до блока 12 не нужны, поскольку owner
остаётся local object.

При success \texttt{commit()} закрывает stream и переименовывает temporary file.
При status early return или exception \texttt{committed\_} остаётся false;
destructor закрывает stream и удаляет temporary path. Release не дублируется.

Это transactional pattern: подготовить новое состояние отдельно, а опубликовать
его после успеха. До commit новый report не виден; при failure temporary state
убирается. Пример не обещает универсальную atomic replace semantics для всех OS
и filesystems: это отдельный contract production utility.

\subsection{Constructor failure}

Если constructor бросает, полностью constructed object не появляется и его
destructor не вызывается как для готового object. Уже успешно constructed
base/member subobjects уничтожаются. Так owner получает сильный
invariant, либо не существовать. В \texttt{TemporaryReport} stream-member уже
создан к моменту проверки; при throw он будет уничтожен автоматически.

\subsection{Destructor и exceptions}

Destructor выполняется и при active unwinding. Если из destructor наружу выходит
второй exception, пока обрабатывается первый, вызывается \texttt{std::terminate}.
Destructors имеют связанные с bases/members implicit \texttt{noexcept} свойства;
здесь достаточно правила: обычный cleanup destructor не должен бросать.

Rollback использует filesystem overload с \texttt{std::error\_code} и не выпускает
cleanup failure. Production code может безопасно фиксировать его в telemetry или
дать отдельную явную \texttt{close}/\texttt{commit}, способную сообщить ошибку до
destructor. Destructor не должен ложно обещать успешное сохранение.

\subsection{Exception-safety guarantees}

Guarantee описывает state после failure:
\textbf{No-throw guarantee.:} Operation обещает не выпускать exception и выполнить свою задачу; обычный кандидат --- простой cleanup destructor; \textbf{Strong guarantee.:} Либо success, либо observable state как до operation; \textbf{Basic guarantee.:} Invariants сохранены и leaks нет, но state мог измениться; \textbf{No guarantee.:} После failure полезного обещания о state нет.

В лаборатории rollback обеспечивает отсутствие нового temporary artifact и
утечки stream на проверяемых failure paths. Prepare-then-commit даёт strong
guarantee относительно появления нового report. Это не доказательство полной
transaction atomicity при любом filesystem failure.

\subsection{Boundary design}

Полезная architecture:
1) low-level operation бросает rich exception об I/O failure; 2) intermediate functions не ловят его без возможности восстановиться; 3) top-level \texttt{main}/command handler ловит по \texttt{const\&}; 4) boundary превращает failure в сообщение и nonzero exit code.

Catch после каждого вызова уничтожил бы отделение failure path. Но ожидаемое
«block incomplete» в course audit не exception: это \texttt{BlockResult}, который
отчёт обязан показать. Выбран hybrid: domain-negative results являются data;
невозможность выполнить audit переносится exception; CLI --- основная catch
boundary. Status mode оставлен рядом для сравнения того же domain logic.

\subsection{Exceptions или status: честное сравнение}

\begin{center}
\begin{tabularx}{\textwidth}{L{0.23\textwidth}YY}
\toprule
Критерий & Status/result & Exception \\
\midrule
Ожидаемая альтернатива & Явна в return type & Менее ясна как частая ветвь \\
Propagation & Повторные checks & До matching handler \\
Normal algorithm & Смешан с transport & Отделён от failure path \\
Control flow & Локально виден & Non-local, требует дисциплины \\
Cleanup & RAII работает & RAII и unwinding вызывают destructors \\
Rich context & Нужен error type & Переносит typed object \\
Цена & Checks на normal path & Throw path дорог; возможны tables/code size \\
\bottomrule
\end{tabularx}
\end{center}

Стоимость non-throwing path зависит от compiler и ABI. Реальный throw обычно
существенно дороже return, поэтому exceptions не используют как частый loop
control. Неверны обе крайности: «exceptions всегда медленные» и «бесплатны».
Destructor work RAII выполняется при окончании lifetime независимо от транспорта.
Стабильный benchmark для этого design conclusion не нужен: BENCHMARKS=N/A.

\subsection{Standard library использует разные модели}

\texttt{vector::at} бросает \texttt{std::out\_of\_range}. Filesystem API часто
предлагает throwing overload и overload со \texttt{std::error\_code}. Streams по
умолчанию имеют state flags и могут быть configured бросать по exception mask.
Не все failures standard library обязательно throw: читают contract операции.

\subsection{Эксперимент 1: error\_model\_lab}

Обе версии решают одну задачу: low-level читает config, middle преобразует его и
пишет report, top выдаёт user-facing result. Status version проверяет
\texttt{Status} на трёх уровнях. Exception version бросает внизу, middle обычно
не ловит, а \texttt{main} ловит по \texttt{const\&}.

Modes детерминированно вызывают failure после acquisition, trace unwinding и
handler с \texttt{throw;}. PASS: success создаёт final report; failure даёт
ожидаемый code; temporary file удалён; trace показывает reverse destruction.

\begin{shellcode}
cd examples/10_exceptions_and_raii/error_model_lab
cmake -S . -B build -DCMAKE_BUILD_TYPE=Debug
cmake --build build
ctest --test-dir build --output-on-failure
./build/error_model_lab unwind build
./build/error_model_lab rethrow build
\end{shellcode}

\subsection{Эксперимент 2: course\_audit\_error\_boundary}

Продолжение \texttt{course\_audit\_patterns} и
\texttt{generic\_course\_audit} сохраняет domain model, но делает error policy
явной. Utility проверяет 15 pairs. Заготовки 11--15 дают нормальные
\texttt{INCOMPLETE}; synthetic I/O failure --- operational failure.

Exception mode не размазывает catch по \texttt{line\_count}, \texttt{inspect} и
\texttt{audit\_exception}; \texttt{main} переводит failure в сообщение и code 2.
Status mode преобразует тот же failure в \texttt{Status}. Основная project policy
--- hybrid, потому что она не путает future content state с I/O failure.

\begin{shellcode}
cd examples/10_exceptions_and_raii/course_audit_error_boundary
cmake -S . -B build -DCMAKE_BUILD_TYPE=Debug
cmake --build build
ctest --test-dir build --output-on-failure
./build/course_audit_error_boundary --root ../../.. --transport exception
\end{shellcode}

PASS: обе transport versions проверяют 15 blocks; expected incomplete не даёт
failure exit; synthetic I/O failure даёт nonzero code и boundary diagnostic.
Sanitizers не применялись: manual memory ownership и намеренного UB здесь нет.

\subsection{Проверка понимания}

После блока нужно уметь:
1) отличить domain result, precondition violation и operation failure; 2) объяснить \texttt{throw}, \texttt{try}, matching \texttt{catch} и поиск наружу; 3) мотивировать catch standard exception по \texttt{const\&}; 4) отличить \texttt{throw;} от \texttt{throw error;}; 5) по trace восстановить destruction при unwinding; 6) определить RAII без smart pointers и назвать разные resources; 7) объяснить cleanup при return, normal exit и exception; 8) описать constructor failure и риск throwing destructor; 9) различить no-throw, strong, basic и no guarantee; 10) выбрать status или exception по contract, frequency и boundary.

\subsection{Источники для сверки}

Throwing, handler matching, rethrow, termination, construction/destruction и
exception safety сверены с C++17 working draft N4659, прежде всего clauses 15 и
18, lifetime --- также с clause 6. RAII сопоставлен с локальными MIT 6.096
materials по memory management и блоками 6, 8, 9. Cost claims ограничены
implementation-dependent формулировками C++ Performance TR; build согласован с
target-based Modern CMake. Старые материалы не использованы как modern style.
